% from ega0-19-fr.tex
\oldpage[0\textsubscript{IV-1}]{107}

$W(k)/pW(k)=k$, $pW(k)$ est l'idéal maximal de $W(k)$, ce qui entraîne que ce
dernier anneau est un \emph{anneau de valuation discrète complet} (Bourbaki,
\emph{Alg. comm.}, chap.~VI, §~3, n\textsuperscript{o}~5, prop.~9). Par
(19.7.1.1) on voit en outre (puisque $k$ est séparable sur $P$, donc une
$P$-algèbre formellement lisse (19.6.1)) que $W(k)$ est une $W(P)$-algèbre
\emph{formellement lisse}. Le $W(P)$-homomorphisme continu $W(k)\to k$ se
factorise donc en $W(k)\to A\to k$ (19.3.11), ce qui permet de considérer $A$
comme une $W(k)$-algèbre topologique. Appliquant maintenant le cas I) à $B_0$
considéré comme $P$-algèbre et à $W(P)$, on voit qu'il existe une
$W(P)$-algèbre $B_P$ qui est un anneau local noethérien complet, un
$W(P)$-module plat, et telle que $B_P\otimes_{W(P)}P$ soit $P$-isomorphe à
$B_0$. Utilisant de nouveau le fait que $W(k)$ est une $W(P)$-algèbre
formellement lisse, on voit par (19.3.11) que l'homomorphisme composé
$W(k)\to k\to B_0$ se factorise en $W(k)\to B_P\to B_0$~; en outre, comme
$k=W(k)/pW(k)$, on a
$B_P\otimes_{W(k)}k=B_P/pB_P=B_P\otimes_{W(P)}P=B_0$. Montrons que $B_P$ est
un $W(k)$-module \emph{plat}~; comme $W(k)$ est un anneau de valuation discrète
dont $p$ est l'uniformisante, il suffit de vérifier que $B_P$ est un
$W(k)$-module \emph{sans torsion} ($\mathbf 0_{\mathrm I}$, 6.3.4), ou encore
que $p$ est un élément $B_P$-régulier, ce qui résulte de ce que $B_P$ est un
$W(P)$-module plat ($\mathbf 0_{\mathrm I}$, 6.3.4).

Posons maintenant
\[
  B=B_P\widehat\otimes_{W(k)}A
\]
et notons que le corps résiduel de $A$ étant égal à celui de $W(k)$, est
\emph{a fortiori} un $W(k)$-module de type fini. Il résulte donc tout d'abord
de (19.7.1.2) que $B$ est un anneau semi-local noethérien complet,
$\mathfrak mB$ étant contenu dans le radical de $B$~; en outre
$B/\mathfrak mB$ est $k$-isomorphe à
$B_P\otimes_{W(k)}k=B_0$, donc $B$ est en fait un anneau local. Comme $B_P$
est un $W(k)$-module plat, (19.7.1.2) montre enfin que $B$ est un $A$-module
plat. C.Q.F.D.
\end{proof}

\begin{lemma}[19.7.1.4]
\label{0.19.7.1.4-fr}
Soient $A$ un anneau, $\mathfrak J$ un idéal de $A$, $M$, $N$ deux $A$-modules
séparés pour la topologie $\mathfrak J$-préadique. On suppose en outre que $N$
est complet pour la topologie $\mathfrak J$-préadique et que $M$ soit un
$A$-module plat. Soit $u:N\to M$ un $A$-homomorphisme~; si
$u\otimes1:N\otimes_A(A/\mathfrak J)\to M\otimes_A(A/\mathfrak J)$ est
bijectif, alors $u$ est bijectif.
\end{lemma}

\begin{proof}
Les modules gradués associés étant pris relatifs aux filtrations
$\mathfrak J$-préadiques, il résulte des hypothèses sur $M$ et $N$ relativement
aux topologies $\mathfrak J$-préadiques qu'il suffit de prouver que
$\gr(u):\gr_\bullet(N)\to\gr_\bullet(M)$ est \emph{bijectif} (Bourbaki,
\emph{Alg. comm.}, chap.~III, §~2, n\textsuperscript{o}~8, cor.~3 du th.~1).
Or, on a un diagramme commutatif
\[
\xymatrix{
  \gr_0(N)\otimes_{A_0}\gr_\bullet(A)
    \ar[r]^-{\gr_0(u)\otimes1}\ar[d]_{\varphi_N} &
  \gr_0(M)\otimes_{A_0}\gr_\bullet(A)\ar[d]^{\varphi_M} \\
  \gr_\bullet(N)\ar[r]_-{\gr(u)} & \gr_\bullet(M)
}
\]

